\documentclass[envcountsect,usenames,dvipsnames,aspectratio=43]{beamer}


\usepackage{pythontex} 
\usefonttheme[onlymath]{serif}
%\usepackage[latin1]{inputenc}
\usepackage[utf8x]{inputenc}
\usepackage[brazilian]{babel}
\usepackage{amsthm,amssymb}
\usepackage{graphicx}
\usepackage{wrapfig}
\usepackage{xcolor}
\usepackage{multicol}
\usepackage{syntonly}


\usepackage{tikz}
\usetikzlibrary{decorations.pathmorphing,patterns}


%%%%%%%%%%%%%%%%%%%%%%%%%%%%%%%%%%%%%%%%%%%%%
% gráfico do coseno nas vibrações mecânicas
\usepackage{pgfplots}
\pgfplotsset{compat=newest}
%%%%%%%%%%%%%%%%%%%%%%%%%%%%%%%%%%%%%%%%%%%%%%

\usepackage[labelformat=empty]{caption}
\setbeamertemplate{caption}[numbered]{}



\usetheme{Madrid}
\usecolortheme{beaver}

\setbeamertemplate{theorems}[numbered]


%%%%%%%%%%%%%%%%%%%%%%%%%%%%%%%%%%%%%%%%%%%%%%%%%%%%%%%%%%%%%%%%%%%%%%%%%%%%%%%%%%%%%%%


\setbeamertemplate{enumerate items}[circle]

\newtheorem{nada}{Nada}

\newtheorem{defin}[nada]{Defini\c c\~ao}

\newtheorem{prop}[nada]{Proposi\c c\~ao}

\newtheorem{corol}[nada]{Corol\'ario}

\newtheorem{teo}[nada]{Teorema}

\newtheorem{lema}[nada]{Lema}

%\newtheorem{exe}{Exemplo}

\newtheorem{outline}{\timesbold outline rem proof}

\newtheorem{obs}[nada]{Observação}
\newtheorem{afir}{Afirmação}

%%%%%%%%%%%%%%%%%%%%%%%%%%%%%%%%%%%%%%%%%%%%%%%%%%%%%%%%%%%%%%



\makeatletter
\def\th@exercicio{%
	\normalfont % body font
	\def\inserttheoremblockenv{alertblock}  
}
\theoremstyle{exercicio}
\newtheorem*{exer}{\includegraphics[scale=0.06]{w-brainb.png} Exercício}
\makeatother

\newtheorem{casa}{\includegraphics[scale=0.035]{w-homework.png} Para Casa}
\makeatother




\makeatletter
\def\th@something{%
	\normalfont % body font
	\def\inserttheoremblockenv{exampleblock}  
}
\theoremstyle{something}
\newtheorem*{exe}{\includegraphics[scale=0.3]{exemplo.png} Exemplo }
\makeatother


\makeatletter
\def\th@resp{%
	\normalfont % body font
	\def\inserttheoremblockenv{block}  
}
\theoremstyle{resp}
\newtheorem*{resp}{\includegraphics[scale=0.01]{White_check.png} Resposta}
\makeatother

\makeatletter
\def\th@desafio{%
	\normalfont % body font
	\def\inserttheoremblockenv{alertblock}  
}

\theoremstyle{desafio}
\newtheorem*{desafio}{\includegraphics[scale=0.02]{desafio-branco.png} Desafio}
\makeatother




%%%%%%%%%%%%%%%%%%%%%%%%

\newcommand{\cqd}{\hfill \framebox[7pt]{} \mbox{} \medskip}
\newcommand{\dem}{\noindent {\bf Demonstra\c c\~ao:}}
\newcommand{\R}{\mathbb{R}}
\newcommand{\vect}[1]{\overrightarrow{#1}}
\newcommand{\vt}[1]{\overrightarrow{#1}}
\newcommand{\dt}[1]{{\color{blue} #1}}
\newcommand{\ang}{\widehat}
\newcommand{\n}[1]{\|#1\|}
\newcommand{\pe}[2]{\langle #1,#2\rangle}
\newcommand{\sen}{\operatorname{sen}} 
\newcommand{\tg}{\operatorname{tg}}     
\newcommand{\proj}{\operatorname{proj}}
\newcommand{\pv}[2]{\vt{#1}\times \vt{#2}}
\newcommand{\pmt}[3]{[\vt{#1},\vt{#2},\vt{#3}]}
\newcommand{\ex}{\textcolor{structure}{\large{Exemplo\ \ }}}
\newcommand{\dps}{\displaystyle}
\newcommand{\senh}{\operatorname{senh}}
\newcommand{\tgh}{\operatorname{tgh}}
\newcommand{\sech}{\operatorname{sech}}
\newcommand{\fm}{\textordfeminine\ }
\newcommand{\mc}{\textordmasculine}

%%%%%%%%%%%%%%%%%%%%%%%%%%%%%%%%%%%%%%%%%%%%%%%%%%%%%%%%%%%%%%%%%%%%%%%%%%%%%%%%%%%%%%%




\begin{document}

%\setbeamercovered{transparent}
\title[cálculo II: EDO]{ Cálculo II\\
	\normalsize{Equações Diferenciais Ordinárias e Aplicações}}
\author[Prof. Reginaldo Demarque]{Prof. Reginaldo Demarque}



%%%%%%%%%%%%%%%%%%%%%%%%%%%%%%%%%%%%%%%%%%%%%%%%%%%%%%%%%%%%%%%%%%%%%%%%%%%%

%\pgfdeclareimage[]{logo}{uff}
%\logo{\pgfuseimage{logo}}

\logo{\includegraphics[scale=0.03]{UFF_brasao.png}}



\institute[UFF/RCN]{\includegraphics[scale=0.4]{sol-logistico-capa.png}\\
	Universidade Federal Fluminense\\
Instituto de Humanidades e Saúde -- RHS\\
Departamento de Ciências da Natureza -- RCN  \\
Campus de Rio das Ostras
	}
\date[{\color{orange} \today}]{ }


\frame{\titlepage}

%%%%%%%%%%%%%%%%%%%%%%%%%% sumario  %%%%%%%%%%%%%%%%%%%%%%%%%%%%%%%%

\frame{
 \frametitle{Sumário}
 \tableofcontents
}



%\AtBeginSection[]
%{
% \begin{frame}
%
%  \frametitle{Sumário}
%  \tableofcontents[currentsection]
%
% \end{frame}
%}

\section{Introdu��o}


\begin{frame}
\frametitle{Introdu��o}
\begin{scriptsize}

\uncover<1->{Imagine que voc� foi contratado por uma empresa que fabrica caixas sem tampa. Cada caixa � contru�da a partir de uma folha retangular de papel�o medindo 30cm x 50cm. Para se construir uma caixa, um quadrado de lado medindo $x$ cm � retirado de cada canto da folha de papel�o. O problema � determinar o valor de $x$ a fim de que a caixa corespondente tenha maior volume poss�vel. 
\bigskip

\begin{center}
\psset{unit=8mm}
\begin{pspicture}(0,0)(7,4.5)
\psline(1,1)(6,1)(6,4)(1,4)(1,1)
\psline[linestyle=dashed](2,1)(2,2)(1,2)
\psline[linestyle=dashed](5,1)(5,2)(6,2)
\psline[linestyle=dashed](1,3)(2,3)(2,4)
\psline[linestyle=dashed](6,3)(5,3)(5,4)
\psline{|-|}(1,0.7)(6,0.7)
\psline{|-|}(0.7,1)(0.7,4)
\put(-0.5,2.5){30cm}
\put(3,0.3){50cm}
\psline{|-|}(6.3,3)(6.3,4)
\psline{|-|}(5,4.3)(6,4.3)
\put(6.4,3.5){$x$}
\put(5.5,4.4){$x$}
\end{pspicture}
\end{center}
Maximizar a \dt{Fun��o-objetivo:} $V(x)=(30-2x)(50-2x)x$ sujeito �s restri��es: $ 0\leq x\leq 15$.
}

\end{scriptsize}
\end{frame}



\begin{frame}
\frametitle{ }
\begin{scriptsize}

\uncover<1->{Agora imagine que voc� foi contratado por  uma empresa que fornece refei��es a seus funcion�rios, a fim de estabelecer um card�pio que atenda a quantidade di�ria m�nima de vitaminas que um funcion�rio deve consumir  e minimize o custo de compras dos diversos tipos de alimentos. Um nutricionista forneceu uma tabela que especifica a quantidade m�nima de cada tipo de vitamina que deve ser ingerida diariamente por cada funcion�rio:
\bigskip

\begin{center}

\begin{tabular}{|l|c|c|c|c|}
\hline
Tipo de vitamina & A & B & C & E\\ \hline
Quantidade  m�nima (mg) & 3 & 1.1 & 60 & 11\\ \hline
\end{tabular} 
\end{center}}

Ele tamb�m forneceu uma tabela com a quantidade (em mg) de vitaminas em cada 100 gramas de 9 tipos de alimentos diferentes,
\begin{center}
\begin{tabular}{|c|c|c|c|c|c|c|c|c|c|c|}
\hline
 Alimento & A & B & C & E &  & Alimento  & A & B & C & E \\  \hline  
1 & 0.140 & 0.08 & 0 & 1.60 & &  6 & 0.035 & 0.05 & 1 & 0 \\ \hline
 2 & 0.580 & 0 & 0 & 0 & &  7 & 0 & 0.04 & 5 & 0.02 \\ \hline
 3 & 0.150 & 1.30 & 26 & 6.90 & &  8 & 0 & 0.08 & 8 & 0.10 \\ \hline
 4 & 0 & 0.08 & 38 & 0.2 & & 9 & 0 & 0.07 & 25 & 2 \\ \hline
 5 & 26.780 & 0.07 & 13 & 0.07 & &  &  &  &  &  \\ \hline
\end{tabular}
\end{center}

E de uma tabela com o pre�o de 100 gramas de cada tipo de alimento:
\begin{center}
\begin{tabular}{|c|c|c|c|c|c|c|c|c|c|}
\hline
Alimento & 1 & 2& 3 & 4 & 5 & 6 & 7 & 8 & 9\\ \hline
Pre�o & 0.60 & 1.00 & 5.00 & 1.00 & 0.50 & 0.20 & 0.15 & 0.40 & 1.00 \\ \hline

\end{tabular}
\end{center}


\end{scriptsize}
\end{frame}



\begin{frame}
\frametitle{ }
\begin{scriptsize}

\uncover<1->{ Neste tipo de problema, cada alimento � uma \dt{vari�vel de controle}, ou seja, a vari�vel $x_1$ representa a quantidade (em ``pacotes'' de 100g) do alimento 1 que a empresa deve comprar diariamente para cada funcion�rio, a vari�vel $x_2$ representa a quantidade (em ``pacotes'' de 100g) do alimento 2 que a empresa deve comprar diariamente para cada funcion�rio, e assim por diante. Resumintdo, $x_i$ representa a quantidade (em ``pacotes'' de 100g) do alimento $i$ que a empresa deve comprar diariamente para cada funcion�rio, com $i=1,\ldots,9$. }
\bigskip

\uncover<2->{A \dt{fun��o-objetivo} que fornece o custo (que queremos minimizar) associado � compra  dos diversos tipos de alimentos para cada funcion�rio � dada por 
$$C(x_1,x_2,x_3,x_4,x_5,x_6,x_7,x_8,x_9)=0.6x_1+x_2+5x_3+x_4+0.5x_5+0.2x_6+0.15x_7+0.4x_8+x_9.$$ }

\uncover<3->{Consultando as outras tabelas, podemos ver que o problema est� sujeito �s restri��es impostas pelas quantidades m�nimas de cada vitamina, a saber
$$\begin{array}{l}
0.14x_1+0.58x_2+0.15x_3+26.79x_5+0.035x_6\geq 3,\\
\\
0.08x_1+1.3x_3+0.08x_4+0.07x_5+0.05x_6+0.04x_7+0.08x_8+0.07x_9\geq 1.1,\\
\\
26x_3+38x_4+13x_5+x_6+5x_7+8x_8+25x_9\geq 60,\\
\\
1.6x_1+6.9_x3+0.2x_4+0.07x_5+0.02x_7+0.1x_8+2x_9\geq 11,\\
\\
x_i\geq 0, \ \forall i=1,\ldots,9.
\end{array} $$}

\end{scriptsize}
\end{frame}



\begin{frame}
\frametitle{ }

\uncover<1->{Usando as t�cnicas que ser�o vistas neste curso de c�lculo 2 � poss�vel resolver este problema e obter a seguinte solu��o:
$$x_1=6.24,\  x_3=0.112, \ x_5=0.078,\ x_7=11.209$$
$$ x_2=x_4=x_6=x_8=x_9=0$$ }

\end{frame}
\section{EDO's de 1\textordfeminine ordem}


\begin{frame}
\frametitle{ EDO's de 1\textordfeminine ordem}
%\begin{scriptsize}

\uncover<1->{Uma EDO de $1^{\underline{a}}$ ordem é uma equação da forma
$$F(t,y,y')=0.$$ 
Um problema da forma
$$\left\{\begin{array}{l}
F(t,y,y')=0\\
y(t_0)=y_0
\end{array}\right.$$
é dito \dt{problema de valor inicial (PVI)}. Uma \dt{solução geral} de uma EDO de $1^{\underline{a}}$ ordem, é uma família de soluções que dependem de uma constante arbitrária, tal que toda solução particular pode ser obtida desta família por uma escolha apropriada da constante.}

%\uncover<2->{\begin{exe}
%Encontre a solução geral da equação $\frac{dy}{dt}=e^{3t}$ e uma solução do PVI
%$$\left\{\begin{array}{l}
%\dps\frac{dy}{dt}=e^{3t}\\
%\\
%y(1/3)=e/3.
%\end{array}\right.$$
%\end{exe} }

%\end{scriptsize}
\end{frame}


\begin{frame}{ }
	
	\begin{exampleblock}{Modelo Populacional Malthusiano}
		Este tipo de modelo é razoável para descrever populações que tem {\color{cyan}recurso ilimitados para crescimento e ausência de predadores}. 
		\begin{itemize}
			\item {\color{blue}$y(t)$}: número de indivíduos de uma população no instante $t$.
			\item {\color{red}$y'(t)$}: taxa de crescimento de uma população no instante $t$.
			
			\item Supõe-se que a {\color{red} taxa de crescimento} de uma população é proporcional à {\color{blue} população presente} 
			\[{\color{red}y'(t)}=k{\color{blue}y(t)}\]
		\end{itemize}
		
		
		\medskip
		
		Supondo que a população no instante $t=0$ é $y_0$, determine a função $y=y(t)$. Em quanto tempo a população dobra de tamanho?
	\end{exampleblock}
\end{frame}



\begin{frame}
\begin{casa}
\begin{enumerate}
\item Determine uma solução geral para a equação
\[y'(t)=y(t).\]

\item Determine uma solução para o PVI
\begin{equation*}
\begin{cases}
y'=y\\
y(0)=2
\end{cases}
\end{equation*}
\end{enumerate}



\end{casa}
\end{frame}

\begin{frame}{Campos de Direções}
\dt{Campos de Direções} são ferramentas validas no estudo de soluções de equações diferenciais da forma
\[y'(t)=f(t,y),\]
onde $f$ é uma função dada chamada de \dt{função taxa de variação}. Ele é construído desenhando-se em cada ponto de uma malha retangular um segmento de reta  cujo coeficiente angular é valor de $f$ naquele ponto.


\end{frame}

\begin{frame}
\begin{exe}
Campo de direções de $y'=y$
\begin{center}
\includegraphics[scale=0.7]{CD-1.png}
\end{center}
\end{exe}
\end{frame}



\begin{frame}
\frametitle{Equações Separáveis }
%\begin{scriptsize}

\uncover<1->{As EDO's de 1\fm ordem \dt{separáveis} são equções da forma
$$g(y)\frac{dy}{dx}=f(x).$$
Integrando esta equação em relação a $x$ temos que
$$\int g(y)y'dx=\int f(x)dx+C.$$
Fazendo a substituição $u=y(x)$, $du=y'(x)dx$ temos que
$$\int g(u)du=\int f(x)dx+C.$$
Assim, se $G$ é uma primitiva de $g$ temos que 
$$G(y(x))=\int f(x) dx +C$$ }
%\end{scriptsize}
\end{frame}


\begin{frame}{Crescimento Populacional Logístico}
	Vimos que um modelo simples de crescimento populacional é aquele em que se supõe que a taxa de crescimento de uma população \textcolor{blue}{$\frac{dy}{dt}$} é proporcional à população presente \textcolor{blue}{$y(t)$} naquele instante.	O \dt{crescimento logístico}, leva em conta que a população tem um valor máximo sustentável \textcolor{red}{$M$}. Quando a população se aproxima da capacidade máxima, os recursos tornam-se menos abundantes e a taxa de crescimento começa a diminuir. Uma relação simples  que exibe esse comportamento é quando 
	\[\textcolor{blue}{\frac{dy}{dt}}={\color{orange}k}\textcolor{blue}{y}(\textcolor{red}{M}-\textcolor{blue}{y}) \]
	
	%\begin{exe} Biólogos colocaram em um lago 400 peixes e estimaram a capacidade de suporte como 10.000. O número de peixes triplicou no primeiro ano. Encontre uma expressão para o tamanho da população de peixes depois de $t$ anos.
	%\end{exe}	
	%	Usando-se o método de frações parciais, pode-se mostrar que a população é modelada por:
	%	\[{\color{blue}y(t)}=\frac{{\color{red}M}}{1+Ce^{-{\color{orange}k}{\color{red}M}t}},\ t\geq 0.\]
	%	
	
\end{frame}

\begin{frame}{ }


%\begin{minipage}{0.55\textwidth}
\begin{exe}
Considere o problema de crescimento logístico:
	\[\begin{cases}
y'=y(1-y)\\
y(0)=y_0, y_0\geq 0.
\end{cases}
\]
Mostre que a solução geral é dada por:
\[y(t)=\frac{1}{1+Ce^{-t}}, t\in \R.\]
\end{exe}
%\end{minipage}
%\begin{minipage}[c]{0.4\textwidth}
%\begin{figure}
%\includegraphics[scale=0.5]{CD-logistico.png}
%\caption{Campo de direções}
%\end{figure}
%\end{minipage}


\end{frame}

\begin{frame}
\begin{block}{Solução geral}
\[y(t)=\frac{1}{1+Ce^{-t}}, t\in \R.\]
\end{block}
\begin{figure}
\centering
\includegraphics[scale=0.65]{sol-logistico-capa.png}
\caption{Família de Soluções para diversos valores de $C$.}
\end{figure}
\end{frame}


\begin{frame}

\begin{exe}
\[y'=\frac{x^2}{1-y^2}\]
\end{exe}

Solução geral
\[y^3+x^3-3y=C\]

\begin{figure}
\centering
\includegraphics[scale=0.35]{sol-ex-sep.png}
\caption{Família de Soluções para diversos valores de $C$.}
\end{figure}

\end{frame}


\begin{frame}
\frametitle{EDO's Lineares de 1\fm ordem }
%\begin{scriptsize}

As \dt{EDO's lineares de 1\fm ordem} são equações que podem ser escritas da forma
\[\frac{dy}{dt}+{\color{cyan} p(t)}y={\color{cyan} q(t)}.\]

\begin{block}{Técnina do Fator Integrante}
Multiplicamos a equação por {\color{red} fator integrante} função {\color{red} $\mu(t)$ }
\[{\color{red} \mu(t)} y'(t)+\only<1>{{\color{red}\mu(t)} {\color{cyan} p(t)}y}\only<2->{\underbrace{{\color{red}\mu(t)} {\color{cyan} p(t)}}_{{\color{red} \mu'(t) }}y}={\color{red}\mu(t)}{\color{cyan}q(t)}.\]

\uncover<2->{\[\Rightarrow ({\color{red} \mu(t) }y(t))'={\color{red} \mu(t) }{\color{cyan}q(t)}\]
Que pode ser resolvida por integração direta. O {\color{red} fator integrante} ${\color{red} \mu(t) }$ pode ser obtido por
\[{\color{red} \mu(t) }=e^{\int{\color{cyan} p(t)}dt}.\]
}

\end{block}


 
%\uncover<1->{\begin{exe} Obtenha a solução geral da equação
%$\dps y'+\frac{2}{t}y=t$ e esboce um gráfico com as soluções.  
%\end{exe}}


%\end{scriptsize}
\end{frame}



\begin{frame}
\begin{exe}
Determine a solução gerla da equação diferencial
\[y'+\frac{1}{2}y=\frac{1}{2}e^{t/3}.\]
Encontre a solução particular que passa pelo ponto $(0,1)$.
\end{exe}


\begin{center}
\only<1>{\includegraphics[scale=0.6]{CD-2.png}}
\only<2>{\includegraphics[scale=0.6]{CD-3.png}}
\end{center}

\end{frame}

%\begin{frame}{Queda Livre com Resistência do ar}
%	Sejam um corpo de massa {\color{blue} $m$} que está caindo e que sofre uma força de resistência do ar que é proporcional à velocidade do corpo. Adotando-se o referencial positivo para baixo, a velocidade satisfaz a equação:
%	\[mv'+kv=mg\]
%	
%\end{frame}
\begin{frame}{Queda Livre com Resistência do ar}
	Sejam um corpo de massa {\color{blue} $m$} que está caindo e que sofre uma {\color{red}força de resistência do ar que é proporcional à velocidade do corpo}. Adotando-se o referencial positivo para baixo, a velocidade satisfaz a equação:
\[{\color{blue}m}v'+{\color{red}k}v={\color{blue}m}g\]
	
\begin{exe} Um pára-quedista com o seu pára-quedas pesam 70 quilogramas e salta de uma altura de 1400 metros. O pára-quedas abre automaticamente após 5 segundos de queda. Sabe-se que a velocidade limite é de 5 $m/s$. Determine a velocidade que o pára-quedista atinge no momento que o pára-quedas abre. Quanto tempo demora para a velocidade chegar a 5,1 $m/s$. Como varia a altura em função do tempo?
\end{exe}
\end{frame}

\begin{frame}
\begin{casa}
\begin{enumerate}
\item Determine a solução geral da EDO: $y'-2y=4-t$

\item Resolva o PVI 
\[\begin{cases}
ty'+2y=4t^2\\
y(1)=2.
\end{cases}\]

\item Determine uma fórmula geral para as soluções da EDO: $y'+ay=g(t)$, onde $a$ é uma constante.
\end{enumerate}
\end{casa}
\end{frame}




%\begin{frame}
%\frametitle{ }
%
%
%\uncover<1->{\begin{exe} Encontre a solução do PVI
%$$\left\{\begin{array}{l}
%\frac{dy}{dx}=\frac{2x-1}{3y^2-3}\\
%y(1)=0
%\end{array} \right.$$
%\begin{enumerate}
%\item Determine o \dt{intervalo de validade da solução}, ou seja, o maior intervalo contendo $x_0=1$ para o qual a solução $y(x)$ e sua derivada estão definidas.
%
%\item Determine os pontos onde a solução tem um máximo local.
%
%\item Faça um esboço do gráfico da solução.
%\end{enumerate}
%\end{exe} }
%
%\end{frame}




\begin{frame}
\frametitle{ Equações Exatas}
%\begin{scriptsize}

Uma EDO de 1\fm\ ordem 
\[{\color{red}M(x,y)}+{\color{blue}N(x,y)}\frac{dy}{dx}=0,\]
é dita \dt{equaçao diferencial exata} quando existe uma função $\psi$ tal que 
\[\frac{\partial \psi}{\partial x}={\color{red}M} \text{ e } \frac{\partial \psi}{\partial y}={\color{blue}N}.\]
Neste caso, podemos reescrever a EDO da forma:
\[\frac{\partial \psi}{\partial x}+\frac{\partial \psi}{\partial y}\frac{dy}{dx}=0\]
Supondo que $y$ é uma função de $x$, pela regra da cadeia para várias variáveis, temos que
\[\frac{d}{dx}(\psi(x,y(x)))=0,\]
Logo a solução geral da EDO é dada implicitamente por
\[\psi(x,y(x))=C\]

%\end{scriptsize}
\end{frame}

\begin{frame}
\begin{exe}
Resolva a EDO: $2x+y^2+2xyy'=0$.
\end{exe}

\begin{teo}
Suponha que ${\color{red}M}, {\color{blue} N}, \frac{\partial {\color{red}M}}{\partial y},\frac{\partial {\color{blue} N}}{\partial x}$ são contínuas num retângulo $[a,b]\times[c,d]$, então a EDO
\[{\color{red}M(x,y)}+{\color{blue}N(x,y)}\frac{dy}{dx}=0\]
é exata se, e somente se, 
\[\frac{\partial \color{red}M}{\partial y }=\frac{\partial {\color{blue} N}}{\partial x}.\]
\end{teo}
\end{frame}


\begin{frame}
\begin{exe}
Verifique se as EDOs são exatas, em caso afirmativo, determine a solução
\begin{enumerate}
\item $(y\cos x+2xe^y)+(\sen x +x^2e^y-1)y'=0$

\item $(3xy+y^2)+(x^2+xy)y'=0$
\end{enumerate}

\end{exe}
\end{frame}

\begin{frame}{Fator Integrante: EDO exata}
No caso em que a EDO não é exata, podemos tentar torná-la exata através de um fator integrante. Seja 
\[{\color{red}M(x,y)}+{\color{blue}N(x,y)}\frac{dy}{dx}=0,\ \text{ com } {\color{red}M_y}\neq {\color{blue}N_x}.\]
Pode-se obter um fator integrante $\mu$  que torna a EDO exata, da seguinte forma:
\begin{enumerate}
\item $P(x)=\frac{{\color{red}M_y}-{\color{blue}N_x}}{{\color{blue}N}}$ então
$\mu(x)=e^{\int P(x)\,dx}$.


\item $Q(y)=\frac{{\color{blue}N_x}-{\color{red}M_y}}{{\color{red}M}}$ então
$\mu(y)=e^{\int Q(y)\,dy}$.
\end{enumerate}

\begin{exe}
Determine o fator integrante que torne a seguinte EDO exata.
\[(3xy+y^2)+(x^2+xy)y'=0\]
\end{exe}
\end{frame}

\begin{frame}
\begin{casa}
 Encontre a solução geral das EDO's:
\begin{enumerate}

\item $(3xy+y^2)+(x^2+xy)y'=0$

\item $\frac{2y(1+x^2)}{1+2x^2}y'-\frac{2xy^2}{(1+2x^2)^2}=1$
\end{enumerate}
\end{casa}
\end{frame}
%
%\subsection{Aplicações}
%
%\begin{frame}
%\frametitle{Dinâmica Populacional }
%\begin{scriptsize}
%
%\uncover<1->{O modelo mais simples de \dt{crescimento populacional} é aquele em que se supõe que a taxa de crescimento de uma população $\frac{dy}{dt}$ é proporcional à população presente $y(t)$ naquele instante . Um outro modelo, conhecido como \dt{crescimento logístico}, leva em conta que a população tem um valor máximo sustentável $y_M$. Quando a população se aproxima da capacidade máxima, os recursos tornam-se menos abundantes e a taxa de crescimento começa a diminuir. Uma relação simples  que exibe esse comportamento é quando 
%$$\frac{dy}{dt}=ky(y_M-y)$$ 
%
%\begin{exe} Biólogos colocaram em um lago 400 peixes e estimaram a capacidade de suporte como 10.000. O número de peixes triplicou no primeiro ano.
%\begin{enumerate}
%\item Encontre uma expressão para o tamanho da população de peixes depois de $t$ anos.
%
%\item Quanto tempo levará para a população aumentear para 5000?
%\end{enumerate}
%\end{exe}}
%
%\end{scriptsize}
%\end{frame}
%
%
%
%\begin{frame}
%\frametitle{Datação por Carbono 14 }
%\begin{scriptsize}
%
%\uncover<1->{A proporção de carbono 14(radioativo) em relação ao carbono 12 presentes nos seres vivos é constante. Quando um organismo morre a absorção de carbono 14 cessa e a partir de então o carbono 14 vai se transformando em carbono 12 a uma taxa proporcional a quantidade presente. Podemos descrever o problema de encontrar a quantidade de carbono 14 em função do tempo, $y(t)$, como o problema de valor inicial
%$$\left\{\begin{array}{l}
%\frac{dy}{dt}=-ky\\
%y(0)=y_0.
%\end{array} \right.$$ 
%
%\begin{exe} Em um pedaço de madeira é encontrado $1/500$ da quantidade original de carbono 14. Sabe-se que a meia-vida do carbono 14 é 5600 anos, ou seja, que em 5600 anos metade do carbono 14 presente transformou-se em carbono 12. Qual a idade deste pedaço de madeira?
%\end{exe}}
%
%\end{scriptsize}
%\end{frame}
%
%
%
%\begin{frame}
%\frametitle{Lei de Resfriamento de Newton }
%
%
%\uncover<1->{A \dt{lei de resfriamento de Newton} diz que a taxa de variação da temperatura $T(t)$ de um corpo em resfriamento é proporcional à diferença entre a temperatura atual do corpo $T(t)$ e a temperatura constante do meio ambiente $T_m$, ou seja, a temperatura do corpo satisfaz a equação
%$$\frac{dT}{dt}=k(T-T_m)$$ 
%
%\begin{exe}
%
%Supopnha que um corpo humano foi encontrado em um motel a meia noite e sua temperature  era de 30\mc C. A temperatura do quarto (temperatura ambiente) era de 15\mc C (supor constante). Duas horas depois a temperatura do corpo era de  25\mc C. Encontre a hora da morte.
%
% 
%\end{exe}}
%
%
%\end{frame}
%
%
%
%
%\begin{frame}
%\frametitle{Lei de Torricelli }
%
%
%\uncover<1->{ A \dt{Lei de Torricelli} diz que a taxa com que um líquido escoa por um orifício situado a uma profundidade $h$ é proporcional a $\sqrt{h}$, ou seja,
%$$\frac{dV}{dt}=k\sqrt{h}.$$
%
%\begin{exe} Um tambor cilíndrico, de 2 metros de altura e base circular de raio 1 metro, está cheio de água. Se fizermos um furo no fundo e em 30 minutos a água cair pela metade qual a altura $h$ em função do tempo e em quanto tempo o tanque esvazia?
%\end{exe}}
%
%
%\end{frame}
%
%
%\begin{frame}
%\frametitle{Resistência em Fluidos }
%
%
%\uncover<1->{ Um corpo que se desloca em um fluido sofre uma força de resitência que é proporcional a velocidade do corpo. Assim, a velocidade, $v(t)$, de um corpo de massa $m$ é solução da equação
%$$m\frac{dv}{dt}=F-kv,$$
%onde $F$ é a força resultante atuando no corpo.
%
%\begin{exe} Um pára-quedista com o seu pára-quedas pesam 70 quilogramas e salta de uma altura de 1400 metros. O pára-quedas abre automaticamente após 5 segundos de queda. Sabe-se que a velocidade limite é de 5 $m/s$. Determine a velocidade que o pára-quedista atinge no momento que o pára-quedas abre. Quanto tempo demora para a velocidade chegar a 5,1 $m/s$. Como varia a altura em função do tempo?
%\end{exe}}
%
%
%\end{frame}



\begin{frame}
\frametitle{Existência e Unicidade de Soluções  }


\uncover<1->{Ao se trabalhar com equações diferenciais duas perguntas são naturais: Um problema de valor inicial 
\[\begin{cases}
y'(t)=f(t,y)\\
y(t_0)=y_0
\end{cases}\]
sempre tem solução? Se sim essa solução é única? 

\begin{exe} 
%\begin{enumerate}
 O problema 
\[ \begin{cases}
y'=2\sqrt{y}\\
y(0)=0
\end{cases}\] tem infinitas soluções! Para todo $c\geq 0$ são soluções do PVI
$$y(t)=\left\{\begin{array}{ll}
(t-c)^2,& t\geq c\\
0, & t<c.
\end{array}  \right.$$
%\end{enumerate}
\end{exe} }

\end{frame}


\begin{frame}
\frametitle{Teorema de Existência e Unicidade de Soluções para Equações Lineares }
\begin{teo} Considere o problema de valor inicial
\[\begin{cases}
y'+{\color{blue}p(t)}y={\color{blue}q(t)}\\
y({\color{red}t_0})=y_0.
\end{cases} \]
Se ${\color{blue}p(t)}$ e ${\color{blue}q(t)}$ são {\color{blue}contínuas} em um intervalo ${\color{red}I}$ contendo ${\color{red}t_0}$, então o PVI {\color{red}tem uma única solução} em $I$.
\end{teo}



\end{frame}


\begin{frame}[fragile=singleslide]
\begin{pycode} 
import sympy as sp


t = sp.symbols('t')
y = sp.symbols('y', cls=sp.Function)

F=sp.exp(-t)
coef=[t**3,4*t**2]

eq=coef[0]*y(t).diff(t)+coef[1]*y(t)

sol=sp.dsolve((eq-F)/coef[0], y(t) )

pvi=sp.dsolve((eq-F)/coef[0], y(t), ics={y(-1): 0})
\end{pycode} 

\begin{exe}
 Considere o PVI:
\[
\begin{cases}
\displaystyle\py{sp.latex(eq)}=\py{sp.latex(F)},\\[1em]
y(-1)=0.
\end{cases}
\]

\begin{enumerate}

\item  Determine o maior intervalo onde a solução do PVI existe.

\item  Encontre a solução do PVI.
\end{enumerate}
\end{exe}
\end{frame}

\begin{frame}
\frametitle{Teorema de Existência e Unicidade de Soluções Geral}


\uncover<1->{ \begin{teo} Considere o problema de valor inicial
\[\begin{cases}
y'(t)={\color{blue}f(t,y)}\\
y({\color{red}t_0})={\color{red}y_0}
\end{cases}\]
Se ${\color{blue}f(t,y)}$ e ${\color{blue}\dps\frac{\partial f}{\partial y}}$ são {\color{blue}contínuas} em um retângulo
\[R=\{(t,x)\in \R^2; a<t<b, c<y<d \}\]
contendo ${\color{red}(t_0,y_0)}$, então o PVI {\color{red}tem uma única solução} em um intervalo contendo $t_0$.
\end{teo}}

\end{frame}





\begin{frame}
\frametitle{ }


\uncover<1->{\begin{exe}
\begin{enumerate}
 \item A única solução do problema
\[\begin{cases}
y'(t)=-y^2\\
y(0)=1
\end{cases}\]
é {\color{blue}$y=\dps\frac{1}{t+1}$} definida no intervalo {\color{red}$(-\infty,-1)$}. Note que não existe uma solução definida em toda a reta!


\item O problema
\[\begin{cases}
y'=\sen(ty)+y^2\\
y(0)=1
\end{cases}\]
tem solução?

\end{enumerate}
\end{exe} }


\end{frame}



\subsection*{Números Complexos}


\begin{frame}{Uma breve introdução aos números complexos\footnotemark}
\footnotetext[1]{Veja o apêndice do livro James Stewart, Cálculo Volume 1.}

 O conjunto dos \dt{números complexos}, denotado por  $\mathbb{C}$, são formados pelos elementos da forma
\[z=a+bi,\ a,b\in \R, \text{ com }  i^2=-1,\]
onde estão definidas as operações de adição e multiplicação:
\begin{itemize}
\item $(a+bi)+(c+di)=(a+c)+(b+d)i$
\item $(a+bi)\cdot(c+di)=(ac-bd)+(ad+bc)i$
\end{itemize}

\begin{exe}
 Calcule $(-1+3i)(2-5i)$
\end{exe}

\end{frame}


\begin{frame}
O  \dt{módulo} de $z$ é definido por 
\[|z|:=\sqrt{a^2+b^2}.\] 

O \dt{Complexo Conjugado} de $z=a+bi$ é 
\[\bar{z}=a-bi.\]


\begin{block}{Propriedades}
\begin{enumerate}
\item $z\bar{z}=|z|^2$.
\item $\overline{z+w}=\bar{z}+\bar{w}$,\ \ \ \ $\overline{zw}=\bar{z}\bar{w}$\ \ \ \ e\ \ \ \ $\overline{z^n}=\bar{z}^n$

\item Se $z$ é raíz de um polinômio, então $\bar{z}$ também o é.
\end{enumerate}
\end{block}



\end{frame}

\begin{frame}
\begin{exe}
\begin{enumerate}
\item Expresse o número $\frac{-1+3i}{2+5i}$ na forma $a+bi$.

\item Encontre as raízes da equação $x^2+x+1$
\end{enumerate}
\end{exe}

\begin{teo}[Teorema Fundamental da Álgebra]
Toda equação polinômial 
\[a_nx^n+a_{n-1}x^{n-1}+\cdots a_1x+a_0=0,\]
onde $a_k\in \R$ e $n\geq 1$, tem exatamente $n$ soluções no conjunto dos números complexos.
\end{teo}
\end{frame}

\begin{frame}{Forma Polar}

Qualquer número complexo $z$ pode ser escrito na {\color{blue} forma polar}
\[z=r(\cos(\theta)+i\sen(\theta)),\]
onde $r=|z|$ e $\tan(\theta)=b/a$. O ângulo $\theta$ é chamado \dt{argumento} de $z$.

\begin{exe}
Escreva $z=1+i$ e $w=i$ na forma polar.
\end{exe}


\end{frame}

\begin{frame}{Propriedades}
Sejam ${\color{blue}z_1}={\color{blue}r_1}(\cos({\color{blue}\theta_1})
+i\sen({\color{blue}\theta_1}))$ ${\color{red}z_2}={\color{red}r_2}(\cos({\color{red}\theta_2})
+i\sen({\color{red}\theta_2}))$. Então
\[{\color{blue}z_1}{\color{red}z_2}={\color{blue}r_1}{\color{red}r_2}
\left(\cos({\color{blue}\theta_1}+{\color{red}\theta_2})
+i\sen({\color{blue}\theta_1}+{\color{red}\theta_2})\right),\]
\[\frac{{\color{blue}z_1}}{{\color{red}z_2}}
=\frac{{\color{blue}r_1}}{{\color{red}r_2}}
\left(\cos({\color{blue}\theta_1}-{\color{red}\theta_2})
+i\sen({\color{blue}\theta_1}-{\color{red}\theta_2})\right),\]
\begin{block}{Fórmula de De Moivre}
Se $z=r(\cos(\theta)+i\sen(\theta))$, então
\[z^n=r^n(\cos(n\theta)+i\sen(n\theta)), \ \forall n\in \mathbb{N}.\]
\end{block}
\end{frame}


\begin{frame}
\begin{exe}
\begin{enumerate}
\item Ache o produto dos números $z=1+i$ e $w=\sqrt{3}-i$ na forma polar.

\item Ache $\left(\frac{1}{2}+\frac{1}{2}i\right)^{10}$.

\item O que significa geometricamente multiplicar um número complexo $z$ por $i$?

\item Calcule $\sqrt[3]{1}$ em  $\mathbb{C}$, isto é, os valores de $z\in \mathbb{C}$ tais que $z^3=1$.
\end{enumerate}
\end{exe}
\end{frame}
\section{EDO's de 2\textordfeminine ordem lineares }




\begin{frame}
\frametitle{ EDO's de 2\fm ordem lineares}

As \dt{EDO's de 2\fm ordem linear } são equações que podem ser escritas na forma
\[y''+{\color{blue}p(t)}y'+{\color{blue}q(t)}y={\color{red}f(t)}.\]

Uma EDO de 2\fm ordem linear é dita \dt{homogênea} se ela pode ser escrita como
\begin{equation}\label{eq_hom_2}
y''+{\color{blue}p(t)}y'+{\color{blue}q(t)}y=0.
\end{equation}

\begin{exe}
\begin{enumerate}
\item EDO Linear de 2\fm ordem não-homogênea: $y''+4y=e^t\sen t$

\item EDO Linear de 2\fm ordem homogênea: $x^2y''+xy'+(x^2-1)y=0$

\item EDO não-Linear de 2\fm ordem: $yy''+y'=0$
\end{enumerate}
\end{exe}
\end{frame}



\begin{frame}
\frametitle{ }
\begin{teo}[Teorema de Existência e Unicidade das Soluções]
Considere o PVI
\[
\begin{cases}
y''+{\color{blue}p(t)}y'+{\color{blue}q(t)}y={\color{blue}f(t)}\\
y({\color{red}t_0})= {\color{red}y_0},\  y'({\color{red}t_0})={\color{red}y_1}.
\end{cases}
\]

Se {\color{blue}$p(t)$, $q(t)$ e $f(t)$} são funções contínuas em um intervalo ${\color{red}I}$ contendo ${\color{red}t_0}$, então o PVI tem {\color{blue} uma única solução definida neste intervalo}.
\end{teo}

\begin{exe} Encontre o maior intervalo no qual a solução do PVI certamente existe.
\[\begin{cases}
(t^2-3t)y''+ty'-(t+3)y=0\\
y(1)= 2,\  y'(1)=1.
\end{cases} 
\]


\end{exe}
\end{frame}

\subsection*{EDO's Homogêneas}

\begin{frame}{EDOs Lineares Homogêneas}

\begin{block}{Princípio da Superposição de Soluções}
Para \dt{EDO's lineares homogêneas}, se $y_1(t)$ e $y_2(t)$ são soluções da equação definidas em um mesmo intervalo, então $$y(t)=c_1y_1(t)+c_2y_2(t)$$
também o é, para quaisquer constantes $c_1$ e $c_2$. 
\end{block}


\begin{exe}Mostre que $y_1(x)=x$ e $y_2(x)=x^3$ são soluções da EDO mas não são soluções do PVI.
\[\begin{cases}
 x^2y''-3xy'+3y=0\\
y(1)= 2,\  y'(1)=1.
\end{cases}
\]


\end{exe}


\end{frame}

\begin{frame}

\begin{casa}
Mostre que:

\begin{enumerate}
\item As funções  $y_1=\cos x$ e $y_2=\sen x$ são soluções da EDO
$$y''+y=0$$

\item As funções $y_1=1+\cos x$ e $y_2=1+\sen x$ são soluções da EDO
$$y''+y=1,$$
mas $y_1+y_2$ não é. 

\item As funções $y_1=x^2$ e $y_2=1$ são soluções da EDO
$$y''y-xy'=0,$$
mas $y_1+y_2$ não é.
\end{enumerate}
\end{casa}
\end{frame}


%
%
%\begin{frame}
%\frametitle{  }
%\begin{scriptsize}
%
%\uncover<1->{ No último exemplo conhecíamos duas soluções da EDO, mas não uma   solução do PVI. Será possível determinar uma solução do PVI a partir dessas duas? E uma solução geral da EDO?}
%\bigskip
%
%\uncover<2->{Seja $\mathcal{S}$ o conjunto de todas as soluções da EDO do último PVI definidas no intervalo $(0,+\infty)$. Da superposição de soluções vemos que $\mathcal{S}$ é um \dt{subespaço vetorial} do espaço vetorial das funções contínuas definida no intervalo $(0,+\infty)$.}
%\bigskip
%
%\uncover<3-> {Defina a transformação 
%$$T:  \mathcal{S} \to  \R^2,\ \  \mbox{ por } \ \  T(y)=(y(1),y'(1)).  $$ 
%Note que
%\begin{itemize}
%\item $T$ é linear ( Exercício);
%\item $T$ é um isomorfismo.
%\end{itemize}
%Portanto, $\dim\mathcal{S}=2$. Com isso, duas soluções LI da EDO formam uma base de $\mathcal{S}$. Como $T(y_1)=(1,1)$ e $T(y_2)=(1,3)$ temos que $y_1$ e $y_2$ são LI, portanto qualquer solução da EDO é da forma
%$$ y(x)=c_1y_1(x)+c_2y_2(x), \forall x\in (0,+\infty)$$
%onde $c_1,c_2\in \R$.}
%
%
%\end{scriptsize}
%\end{frame}
%
%




\begin{frame}
\frametitle{ Soluções Fundamentais }
 No último exemplo vimos que $y_1(x)=x$ e $y_2(x)=x^3$, $\forall\ x\in (0,+\infty)$ são soluções da EDO $x^2y''-3xy'+3y=0$ mas não do PVI
\[\begin{cases}
 x^2y''-3xy'+3y=0\\
y(1)= 2,\  y'(1)=1.
\end{cases}
\]

Será possível determinar uma solução do PVI a partir dessas duas? E uma solução geral da EDO?

\end{frame}
%
%
%\begin{frame}
%\frametitle{}
%\begin{scriptsize}
%
%\uncover<1->{ Com isso, uma solução geral da EDO é da forma
%\begin{equation}\label{edo1}
%y(x)=c_1x+c_2x^3, \forall x\in (0,+\infty).
%\end{equation}
%Para obtermos uma solução do PVI devemos determinar $c_1$ e $c_2$  tal que $T(y)=(2,1)$, ou seja,
%$$(2,1)=T(y)=c_1(1,1)+c_2(1,3)\Rightarrow \left\{\begin{array}{l}
%c_1+c_2=2\\
%c_1+3c_2=1.
%\end{array}\right.\Rightarrow c_1=\frac{5}{2} \mbox{ e } c_2=-\frac{1}{2}.$$
%Logo, a solução do PVI é 
%$$y(x)=\frac{5x}{2}-\frac{x^3}{2}, \forall x\in (0,+\infty).$$  }
%
%
%\end{scriptsize}
%\end{frame}


%\begin{frame}
%\frametitle{ Soluções Fundamentais}
%\begin{scriptsize}
%
%
%
%\uncover<1->{Considere o PVI
%$$\left\{\begin{array}{l}
%y''+p(t)y'+q(t)y=0\\
%y(t_0)= a,\  y'(t_0)=b.
%\end{array} \right.$$
%Se $p(t)$ e $q(t)$ são contínuas, então o procedimento do exemplo anterior pode ser aplicado.}
%\bigskip
%
%\uncover<2->{Sejam $y_1$ e $y_2$ duas soluções quaisquer da EDO. Dado $\tilde{y}\in I$ qualquer, como $T_{\tilde{t}}$ é um isomorfismo, sabemos que $y_1$ e $y_2$ são LI se, e somente se, $T_{\tilde{t}}(y_1)=(y_1(\tilde{t}),y'_1(\tilde{t}))$ e $T_{\tilde{t}}(y_2)=(y_2(\tilde{t}),y'_2(\tilde{t}))$ são LI, ou ainda, se, e somente se, 
%$$\det\left(\begin{array}{cc}
%y_1(\tilde{t}) & y'_1(\tilde{t})\\
%y_2(\tilde{t}) & y'_2(\tilde{t})
%\end{array}\right)\neq 0$$}
%
%\end{scriptsize}
%\end{frame}
%
%
%\begin{frame}
%\frametitle{ }
%\begin{scriptsize}
%
%\uncover<1->{Se $y_1$ e $y_2$ são duas soluções LI, então a solução do PVI é da forma
%$y=c_1y_1+c_2y_2$. Das condições iniciais temos que
%$$
%\left\{\begin{array}{l}
% c_1y_1(t_0)+c_2y_2(t_0)=a\\
% c_1y'_1(t_0)+c_2y'_2(t_0)=b
%\end{array}\right.
%$$}
%
%\uncover<1->{Como $y_1$ e $y_2$ são  duas soluções LI, então  
%$$\det\left(\begin{array}{cc}
%y_1(t_0) & y_2(t_0)\\
%y'_1(t_0) & y'_2(t_0)
%\end{array}\right)\neq 0,$$
%Portanto o sistema linear possui única solução! }
%\bigskip
%
%\uncover<2->{ O último determinante é dito \dt{wronskiano} das funções $y_1$ e $y_2$ em $t_0$ e é denotado por $W[y_1,y_2](t_0)$. Se $y_1$ e $y_2$ são tais que $W[y_1,y_2](t)\neq 0$ para algum $t\in I$ dizemos que elas são  \dt{soluções fundamentais}.}
%
%
%\end{scriptsize}
%\end{frame}




\begin{frame}
\frametitle{Wronskiano}
Considere o PVI
\[
\begin{cases}
y''+{\color{blue}p(t)}y'+{\color{blue}q(t)}y=0\\
y({\color{red}t_0})= {\color{red}y_0},\  y'({\color{red}t_0})={\color{red}y_1}.
\end{cases}
\]
Se ${\color{blue}p(t)}$ e ${\color{blue}q(t)}$ são contínuas, então o procedimento do exemplo anterior pode ser aplicado. Dados {\color{cyan} $y_1$} e {\color{cyan} $y_2$} duas da EDO, então o PVI terá solução desde que 
\[W[{\color{cyan} y_1},{\color{cyan} y_2}]({\color{red}t_0})=
\det\left( \begin{array}{cc}
 {\color{cyan} y_1}({\color{red}t_0}) & {\color{cyan} y_2}({\color{red}t_0})  \\
{\color{cyan} y_1^\prime}({\color{red}t_0}) & {\color{cyan} y_2^\prime}({\color{red}t_0})
\end{array}
\right).\]
$W[{\color{cyan} y_1},{\color{cyan} y_2}]({\color{red}t_0})$ é chamado de {\color{blue} Wronskiano}. 


\end{frame}



\begin{frame}
\begin{teo}
Se {\color{cyan} $y_1$} e {\color{cyan} $y_2$} são soluções da EDO 
\[y''+{\color{blue}p(t)}y'+{\color{blue}q(t)}y=0,\]
e se existe ${\color{red}t_0}$ tais que $W[{\color{cyan} y_1},{\color{cyan} y_2}]({\color{red}t_0})\neq 0$, então a família de funções
\[y=c_1{\color{cyan} y_1}+c_2{\color{cyan} y_2},\]
incluem todas as soluções da EDO, chamada \dt{solução geral} da EDO. Neste caso, {\color{cyan} $y_1$} e {\color{cyan} $y_2$} são ditas \dt{soluções fundamentais}.
\end{teo}

\begin{exe}
Mostre que $y_1=t^{1/2}$ e $y_2=t^{-1}$ são soluções fundamentais da EDO
\[2t^2y''+3ty'-y=0,\ t>0\]
e determine a solução geral.
\end{exe}
\end{frame}



%
%
%
%%
%%
%\begin{frame}{Vibrações Mecânicas}
%\begin{tikzpicture}[black!75]
%		
%		% Supporting structure
%		\fill [pattern = north west lines] (-1.5,0) rectangle ++(3,.2);
%		\draw[thick] (-1.5,0) -- ++(3,0);
%		
%		% Spring + Arrows
%		\draw[] (0,0) -- ++(0,-0.25);
%		\draw[decoration={aspect=0.3, segment length=1.2mm, amplitude=2mm,coil},decorate] (0,-0.25) -- ++(0,-2.25) node[midway,right=0.25cm,black]{{\color{red}$k$}}; 
%		\draw[] (0,-2.5) -- ++(0,-0.3) node[coordinate](c1){};
%		
%		\begin{scope}[xshift=4cm]
%			% Supporting structure
%			\fill [pattern = north west lines] (-1.5,0) rectangle ++(3,.2);
%			\draw[thick] (-1.5,0) -- ++(3,0);
%			
%			% Spring + Arrows
%			\draw[] (0,0) -- ++(0,-0.25);
%			\draw[decoration={aspect=0.3, segment length=1.4mm, amplitude=2mm,coil},decorate] (0,-0.25) -- ++(0,-2.75) node[midway,right=0.25cm,black]{{\color{red}$k$}}; 
%			\draw[] (0,-3) -- ++(0,-0.3)node[coordinate](c2){} node[draw,fill=blue!70,minimum width=1cm,minimum height=0.5cm,anchor=north,label=east:{\color{blue}$m$}](M){};
%		\end{scope}
%		
%		\begin{scope}[xshift=8cm]
%			% Supporting structure
%			\fill [pattern = north west lines] (-1.5,0) rectangle ++(3,.2);
%			\draw[thick] (-1.5,0) -- ++(3,0);
%			
%			% Spring + Arrows
%			\draw[] (0,0) -- ++(0,-0.25);
%			\draw[decoration={aspect=0.3, segment length=1.5mm, amplitude=2mm,coil},decorate] (0,-0.25) -- ++(0,-3.5) node[midway,right=0.25cm,black]{{\color{red}$k$}}; 
%			\draw[] (0,-3.75) -- ++(0,-0.3)node[coordinate](c3){} node[draw,fill=blue!70,minimum width=1cm,minimum height=0.5cm,anchor=north,label=east:{\color{blue}$m$}](M){};
%		\end{scope}
%		
%		
%		\draw[dashed,gray] (c1) -- ++(3.75,0)coordinate(c22);
%		\draw[dashed,gray] (c2) -- ++(-1.5,0) coordinate(c12);
%		\draw[latex-latex] (c12)-- (c12|-c1)node[midway,left]{\small $L$};
%
%
%		
%		\draw[dashed,gray] (c22)++(0.5,0) -- ++(3.5,0)coordinate(c33);
%		\draw[dashed,gray] (c3) -- ++(-1.5,0) coordinate(c23);
%
%	    \draw[dashed,gray] (c2)++(.5,0) -- ++(2,0) coordinate(c122);
%		\draw[latex-latex] (c122)-- (c122|-c3)node[midway,left]{\small $y$};
%		
%		
%	\end{tikzpicture}
%
%Um sistema de massa-mola composto de um corpo de massa {\color{blue}$m$} preso a uma mola, com constante elástica {\color{red}$k$}, que está presa ao teto satisfaz a equação diferencial
%\[{\color{blue}m}y''+{\color{red}k}y =0.\]
%\end{frame}
%
%
%\begin{frame}{Pêndulo Simples}
%O movimento de um pêndulo simples de massa $m$ e comprimento $\ell$ é descrito pela função $\theta(t)$ que satisfaz a equação diferencial
%\[\theta''+\frac{g}{\ell}\sen\theta=0.\]
%
%\end{frame}



\begin{frame}{Vibrações Mecânicas Amortecidas}
\begin{tikzpicture}[black!75]
		
		% Supporting structure
		\fill [pattern = north west lines] (-1.5,0) rectangle ++(3,.2);
		\draw[thick] (-1.5,0) -- ++(3,0);
		
		% Spring + Arrows
		\draw[] (0,0) -- ++(0,-0.25);
		\draw[decoration={aspect=0.3, segment length=1.2mm, amplitude=2mm,coil},decorate] (0,-0.25) -- ++(0,-2.25) node[midway,right=0.25cm,black]{{\color{red}$k$}}; 
		\draw[] (0,-2.5) -- ++(0,-0.3) node[coordinate](c1){};
		
		\begin{scope}[xshift=4cm]
			% Supporting structure
			\fill [pattern = north west lines] (-1.5,0) rectangle ++(3,.2);
			\draw[thick] (-1.5,0) -- ++(3,0);
			
			% Spring + Arrows
			\draw[] (0,0) -- ++(0,-0.25);
			\draw[decoration={aspect=0.3, segment length=1.4mm, amplitude=2mm,coil},decorate] (0,-0.25) -- ++(0,-2.75) node[midway,right=0.25cm,black]{{\color{red}$k$}}; 
			\draw[] (0,-3) -- ++(0,-0.3)node[coordinate](c2){} node[draw,fill=blue!70,minimum width=1cm,minimum height=0.5cm,anchor=north,label=east:{\color{blue}$m$}](M){};
		\end{scope}
		
		\begin{scope}[xshift=8cm]
			% Supporting structure
			\fill [pattern = north west lines] (-1.5,0) rectangle ++(3,.2);
			\draw[thick] (-1.5,0) -- ++(3,0);
			
			% Spring + Arrows
			\draw[] (0,0) -- ++(0,-0.25);
			\draw[decoration={aspect=0.3, segment length=1.5mm, amplitude=2mm,coil},decorate] (0,-0.25) -- ++(0,-3.5) node[midway,right=0.25cm,black]{{\color{red}$k$}}; 
			\draw[] (0,-3.75) -- ++(0,-0.3)node[coordinate](c3){} node[draw,fill=blue!70,minimum width=1cm,minimum height=0.5cm,anchor=north,label=east:{\color{blue}$m$}](M){};
		\end{scope}
		
		
		\draw[dashed,gray] (c1) -- ++(3.75,0)coordinate(c22);
		\draw[dashed,gray] (c2) -- ++(-1.5,0) coordinate(c12);
		\draw[latex-latex] (c12)-- (c12|-c1)node[midway,left]{\small $L$};


		
		\draw[dashed,gray] (c22)++(0.5,0) -- ++(3.5,0)coordinate(c33);
		\draw[dashed,gray] (c3) -- ++(-1.5,0) coordinate(c23);

	    \draw[dashed,gray] (c2)++(.5,0) -- ++(2,0) coordinate(c122);
		\draw[latex-latex] (c122)-- (c122|-c3)node[midway,left]{\small $y$};
		
		
	\end{tikzpicture}

Considere um sistema de massa-mola composto de um corpo de massa {\color{blue}$m$} preso a uma mola, com constante elástica {\color{red}$k$}, que está presa ao teto. Se levarmos em conta um amortecimento viscoso proporcional à velocidade do corpo, então o sistema satisfaz a EDO
\[{\color{blue}m}y''+{\color{orange}\gamma} y'+{\color{red}k}y =0,\]
onde ${\color{orange}\gamma}>0$ é a constante de amortecimento.
\end{frame}


\begin{frame}
\frametitle{Equações homogêneas com coeficientes constantes }

Uma EDO linear de 2\fm ordem, homogênea, com coeficientes constantes é uma equação  da forma
\begin{equation}\label{coef_const}
ay''+by'+cy=0,\ a,b,c\in\R,\ a\neq 0.
\end{equation} 
 
\bigskip

Para resolver uma equação do tipo (\ref{coef_const}) vamos nos inspirar no caso de 1\fm ordem. Uma EDO linear homogêna de 1\fm com coeficientes constantes é da forma
\[ay'+by=0,\ a,b\in\R,\ a\neq 0.\]
Sabemos que as soluções para esta equação são $\dps y(t)=ce^{-bt/a}$. Neste caso é natural supor que uma solução da EDO (\ref{coef_const}) seja da forma $y(t)=e^{\lambda t}$ para alguma constante $\lambda$. Daí, substituindo em (\ref{coef_const}) temos que
\[a\lambda^2e^{\lambda t}+b\lambda e^{\lambda t}+ce^{\lambda t}=0\Leftrightarrow a\lambda^2+b\lambda+c=0.\]
A última equação é dita \dt{equação característica.}

\end{frame}


\begin{frame}{Raízes Reais Distintas}
\begin{exe}
 Determinar a solução geral da EDO: $y''+y'-2y=0$.
\end{exe}

\begin{block}{}

Se $\lambda_1$ e $\lambda_2$ são raízes distintas da equação característica, então a solução geral da EDO é:
\[y(t)=c_1e^{\lambda_1t}+c_2e^{\lambda_2t}, t\in \R.\]
\end{block}



\end{frame}

\begin{frame}{Raízes Reais Iguais}




\begin{exe}
 Determinar a solução geral da EDO: $y''+4y'+4y=0$
\end{exe}
\begin{block}{}
Se $\alpha$ é a única raiz da equação característica, então a solução geral da EDO é:
\[y(t)=c_1e^{\alpha t}+c_2te^{\alpha t}, t\in \R.\]


\end{block}

\end{frame}


%\begin{frame}
%
%\begin{exe} Determinar as soluções da equação:
%\begin{enumerate}[a]
%\item $y''+y'-2y=0$
%\item $y''+2y'+y=0$
%\item $y''-2y+2=0$
%\end{enumerate}
%\end{exe}
%\end{frame}

%\begin{frame}
%\frametitle{ }
%\begin{scriptsize}
%
%\uncover<1->{\begin{obs}
%
%\begin{enumerate}[a)]
%
%\item Se um dos coeficientes da EDO, $p(t)$ ou $q(t)$,  não é  contínuo no intervalo $I$, então o Wronskiano pode se anular e as funções serem LI! De fato, se no exemplo anterior procurarmos por soluções definidas em toda a reta, então as funções 
%$$ y_1=x^3\ \  \mbox{ e } \ \ 
%y_2=\left\{ \begin{array}{ll}
%x^3,& x\geq 0\\
%-x^3,& x<0
%\end{array}  \right.$$
%são soluções LI da EDO $y''-\frac{3}{x}y'+\frac{3}{x^2}y=0$ e $W[y_1,y_2](x)=0$ para todo $x\in \R.$ Isso ocorre pois as funções não são contínuas em nenhum intervalo contendo a origem. Entretanto, se restringirmos a intervalos que não contém a origem elas são LD, como por exemplo o intervalo $(-\infty,0)$, neste caso $y_1=-y_2$. Note ainda que $y_3=x$ também é solução da EDO linearmente independente das outras duas, portanto $\dim \mathcal{S}\geq 3$.
%
%\item Entretanto, se $W[y_1,y_2](t_0)\neq 0$ para algum $t_0\in I$, então $y_1$ e $y_2$ são LI em $I$. De fato, se $W[y_1,y_2](t_0)\neq 0$, então $T_{t_0}(y_1)$ e $T_{t_0}(y_2)$ são LI, daí,
%$$\alpha_1 y_1+\alpha_2y_2=0\Rightarrow \alpha_1 T_{t_0}(y_1)+\alpha_2T_{t_0}(y_2)=0\Rightarrow \alpha_1=\alpha_2=0.  $$
% 
%
%\end{enumerate} 
%
%\end{obs}}
%
%
%\end{scriptsize}
%\end{frame}





\begin{frame}
Vamos estudar inicialmente o caso
\[y''+{\color{red}\beta^2}y=0,\]
cujas soluções fundamentais são $y_1=e^{i{\color{red} \beta}t}$ e $y_2=e^{-i{\color{red} \beta}t}$. Mas como podemos escrevê-las na forma padrão  $a+bi$?
\medskip

Primeiramente, note que $u_1(t)=\cos( {\color{red} \beta}t)$ e $u_2(t)=\sen( {\color{red} \beta}t)$ também são soluções fundamentais da EDO. Portanto, 
\[y_1(t)=c_1\cos( {\color{red} \beta}t)+c_2\sen( {\color{red} \beta}t).\]
%Como $y_1(0)=1$ e $y_1'(0)=i{\color{red} \beta}$, daí, temos que $c_1=1$ e $c_2=i$. Logo,
%\[e^{i{\color{red}\beta}t}=\cos( {\color{red} \beta}t)+i\sen( {\color{red} \beta}t).\]
%Em particular, para $t=1$, temos que 
%\[e^{i{\color{red} \beta}}=\cos( {\color{red} \beta})+i\sen( {\color{red} \beta}).\]
\end{frame}

\begin{frame}{A exponencial complexa}
Até agora temos que
\[\begin{cases}
y_1(t)=e^{i{\color{red} \beta}t}\\
y_1(t)=c_1\cos( {\color{red} \beta}t)+c_2\sen( {\color{red} \beta}t).
\end{cases}\]
Como $y_1(0)=1$ e $y_1'(0)=i{\color{red} \beta}$, daí, temos que $c_1=1$ e $c_2=i$. Logo,
\[e^{i{\color{red} \beta}t}=\cos( {\color{red} \beta}t)+i\sen( {\color{red} \beta}t).\]
Em particular, para $t=1$, temos que 
\[e^{i{\color{red} \beta}}=\cos( {\color{red} \beta})+i\sen( {\color{red} \beta}).\]



\end{frame}


\begin{frame}
Curiosamente, quando  ${\color{red} \beta}=\pi$, obtemos {\color {red} a mais bela de todas as equações da matemática}:
\begin{block}{Equação de Euler}
\[e^{i\pi}+1=0.\]
\end{block}



De forma geral, obtemos:
\[e^{{\color{blue}\alpha}+{\color{red} \beta} i}=e^{\color{blue}\alpha} e^{{\color{red} \beta} i}=e^{\color{blue}\alpha}(\cos {\color{red} \beta} +i\sen {\color{red} \beta}),\]
conhecida como fórmula de Euler.
\end{frame}


\begin{frame}{Raízes Complexas}
\begin{exe}
Determine a solução geral da EDO : $y''-4y'+13y=0$.
\end{exe}
\begin{block}{ }
Se $\lambda_1={\color{blue}\alpha}+{\color{red} \beta} i$ e $\lambda_2={\color{blue}\alpha}-{\color{red} \beta} i$, então a solução geral da EDO é:
\[y(t)=e^{{\color{blue}\alpha} t}(c_1\cos({\color{red} \beta} t)+c_2\sen({\color{red} \beta} t)),\ t\in \R.\]
\end{block}
\end{frame}


\section{Vibrações Mecânicas Livres}

\subsection{Vibrações Livres}

\begin{frame}{Vibrações Mecânicas Livres}
Vimos que o modelo para um sistema massa-mola preso no teto em um meio viscoso é:
\[{\color{blue}m}y''+{\color{orange}\gamma} y'+{\color{red}k}y =0,\]
onde {\color{blue}$m$} é a massa, ${\color{red}k}>0$ é a constante elástica e   ${\color{orange}\gamma}>0$ é a constante de amortecimento.
\end{frame}

\begin{frame}{Vibrações livres não-amortecidas}
Quando ${\color{orange}\gamma}=0$, o sistema não tem amortecimento e podemos reescrever a equação:
\[y''+{\color{violet}\omega_0^2}y =0,\]
onde ${\color{violet}\omega_0^2}=\frac{{\color{red}k}}{{\color{blue}m}}$. Com isso a solução geral é:
\[y=A\cos({\color{violet}\omega_0}t)+B\sen({\color{violet}\omega_0}t),\ t\in \R.\]
 A solução pode ser reescrita como:
\[y=R\cos({\color{violet}\omega_0}t-\delta),\]
onde $A=R\cos\delta$, $B=R\sen \delta$. O \dt{período} do movimento é $T=\frac{2\pi}{{\color{violet}\omega_0}}$, a \dt{frequência} é $f=\frac{{\color{violet}\omega_0}}{2\pi}$, a \dt{amplitude} é $R$ e o parâmetro adimensional $\delta$ é chamado de \dt{fase}.

\end{frame}


\begin{frame}
O movimento descrito é chamado \dt{movimento harmônico}.
\begin{center}
\begin{tikzpicture}
  \begin{axis}%
    [clip mode=individual,
	 grid=both,
     minor tick num=4,
     grid style={line width=.5pt, draw=gray!10},
     major grid style={line width=.2pt,draw=gray!50},
     axis lines=middle,
	 yticklabels={,,},
	 xticklabels={,,},
     enlargelimits={abs=0.2}
	]
    \addplot[domain=0:2.6,samples=50,smooth,red,thick] {cos(deg(pi*x-0.4*pi))};
\node[left ] at (0,1) {$R$};
\node[left ] at (0,-1) {$-R$};
\node[below] at (.4,0) {$\frac{\delta}{{\color{violet}\omega_0}}$};
\draw[dashed,cyan,thick] (0.4,0) -- (0.4,1);
\node[below] at (2.4,0) {$\frac{2\pi+\delta}{{\color{violet}\omega_0}}$};
\draw[dashed,cyan,thick] (2.4,0) -- (2.4,1);
\node at (3,0) {$t$};
\node at (0,1.3) {$y$};
\node at (1.5,1.5) {$y=R\cos({\color{violet}\omega_0}t-\delta)$};
  \end{axis}
\end{tikzpicture}
\end{center}

\end{frame}

\begin{frame}
\begin{center}
\includegraphics[scale=0.8]{osc-harm.png}
\end{center}
\end{frame}


\subsection{Vibrações Livres Amortecidas}


\begin{frame}{Vibrações Livres Amortecidas}
Quando o sistema é amortecido temos a equação:
\[{\color{blue}m}y''+{\color{orange}\gamma} y'+{\color{red}k}y =0.\]
Donde temos a equação característica:
\[{\color{blue}m}\lambda^2+{\color{orange}\gamma} \lambda+{\color{red}k} =0,\]
cujas raízes são:
\[\lambda=-\frac{{\color{orange}\gamma}}{2{\color{blue}m}}\pm \frac{\sqrt{{\color{orange}\gamma}^2-4{\color{blue}m} {\color{red}k}}}{2{\color{blue}m}},\]
o que nos fornece 3 casos:

\begin{enumerate}
\item \textbf{Superamortecimento}: \quad ${\color{orange}\gamma}^2>4{\color{blue}m} {\color{red}k}$
\item \textbf{Subamortecimento}: \qquad  ${\color{orange}\gamma}^2<4{\color{blue}m} {\color{red}k}$
\item \textbf{Amortecimento Crítico}: \ ${\color{orange}\gamma}^2=4{\color{blue}m} {\color{red}k}$
\end{enumerate}
\end{frame}


\begin{frame}{Superamortecimento}
 \[y(t)=c_1e^{\lambda_1 t}+c_2e^{\lambda_2 t},\]
onde $\lambda_1,\lambda_2<0$.

\includegraphics[scale=0.45]{osc-super.png}
%\begin{tikzpicture}
%%	\def\m{1}
%%	\def\k{2*\m}
%%	\def\g{3*\m}
%%\def\a{\g/(2*\m)}
%%\def\b{sqrt(\g*\g-4*\m*\k)/(2*\m)}
%%\def\L{-\a+\b}
%%\def\l{-\a-\b}
%  \begin{axis}%
%    [clip mode=individual,
%	 grid=both,
%     minor tick num=4,
%     grid style={line width=.5pt, draw=gray!10},
%     major grid style={line width=.2pt,draw=gray!50},
%     axis lines=middle,
%	 %yticklabels={,,},
%	 %xticklabels={,,},
%     enlargelimits={abs=0.2}
%	]
%    \addplot[domain=0:4,samples=50,smooth,red,thick] {2.7*exp(-\x)-2.1*exp(-2*\x)};
%	\addplot[domain=0:4,samples=50,smooth,cyan,thick] {0.4*exp(-\x)+0.1*exp(-2*\x)};
%	\addplot[domain=0:4,samples=50,smooth,orange,thick] {-3.3*exp(-\x)+3*exp(-2*\x)};
%%\node[left ] at (0,1) {$R$};
%%\node[left ] at (0,-1) {$-R$};
%%\node[below] at (.4,0) {$\frac{\delta}{{\color{violet}\omega_0}}$};
%%\draw[dashed,cyan,thick] (0.4,0) -- (0.4,1);
%%\node[below] at (2.4,0) {$\frac{2\pi+\delta}{{\color{violet}\omega_0}}$};
%%\draw[dashed,cyan,thick] (2.4,0) -- (2.4,1);
%%\node at (3,0) {$t$};
%%\node at (0,1.3) {$y$};
%%\node at (1.5,1.5) {$y=R\cos({\color{violet}\omega_0}t-\delta)$};
%  \end{axis}
%\end{tikzpicture}


\end{frame}


\begin{frame}{Subamortecimento}
A solução é da forma
 \[y(t)=e^{-\frac{{\color{orange}\gamma}t}{2{\color{blue}m}}}(A\cos(\mu t)+B\sen(\mu t)),\]
onde $\mu=\frac{\sqrt{4{\color{blue}m {\color{red}k}}-{\color{orange}\gamma}^2}}{2{\color{blue}m}}$.
Que pode ser reescrita como
 \[y(t)=Re^{-\frac{{\color{orange}\gamma}t}{2{\color{blue}m}}}\cos(\mu t-\delta),\]
onde $A=R\cos\delta$ e $B=R\sen\delta$.

\begin{center}
\includegraphics[scale=0.45]{osc-subarm.png}
\end{center}
\end{frame}




\begin{frame}{Amortecimento crítico}
A solução é da forma
 \[y(t)=(A+Bt)e^{-\frac{{\color{orange}\gamma}t}{2{\color{blue}m}}}.\]


\begin{center}
\includegraphics[scale=0.5]{osc-critico.png}
\end{center}
\end{frame}

\begin{frame}
\begin{casa}
Esboce os gráficos das soluções dos problemas abaixo.
\begin{enumerate}
\item Suponha que uma massa de 4,5 kg estica uma mola x5cm. A massa é deslocada  2,5 cm para baixo e depois colocada em movimento  com uma velocidade inicial de apontando para cima de 30 cm/s. 

\item Uma massa de 20g estica uma mola 5 cm. Suponha que a massa também está presa a um amortecedor viscoso com uma constante de amortecimento de 400 dinas$\cdot$s/cm e que a massa é puxada pra baixo mais 2 cm de depois é solta. 
\end{enumerate}
\end{casa}
\end{frame}

%
%\begin{frame}
%\frametitle{ }
%\begin{scriptsize}
%
%\uncover<1->{\begin{exe} Mostre que $y_1(t)=\cos bt$ e $y_2(t)=\sen bt$ são soluções fundamentais da equação
%$$y''+b^2y=0$$
%
%\end{exe} }
%
%Usando o exemplo anterior podemos estender a função exponencial para os números complexos. Queremos definir $y(t)=e^{zt}$ para números complexos $z=a+bi$ de forma que satisfaça as propriedades
%\begin{align*}
%& e^{(a+ib)t}=e^{at}e^{ibt}\\
%& \frac{d}{dt}(e^{zt})=ze^{zt}.
%\end{align*}
%Note que a função $w(t)=e^{ibt}$ é solução do PVI
% $$\left\{\begin{array}{l}
% y''+b^2y=0 \\
% y(0)=1,\  y'(0)=ib.
% 
%\end{array}  \right.$$
%Assim, do exemplo anterior, como $y_1(t)=\cos bt$ e $y_2(t)=\sen bt$ são soluções fundamentais da EDO, temos que
%$$z(t)=\cos bt+i\sen bt.$$
%
%
%
%\end{scriptsize}
%\end{frame}
%
%
%
%\begin{frame}
%\frametitle{ }
%\begin{scriptsize}
%
%\uncover<1->{ Portanto, das propriedades de exponencial, obtemos que
%$$e^{a+bi}=z(1)=e^a(\cos b+i\sen b),$$
%conhecida como \dt{fórmula de Euler}.}
%
%\end{scriptsize}
%\end{frame}
%
%
%\begin{frame}
%\frametitle{Equações homogêneas com coeficientes constantes }
%\begin{scriptsize}
%
%\uncover<1->{Uma EDO linear de 2\fm ordem, homogênea, com coeficientes constantes é uma equação  da forma
%\begin{equation}\label{coef_const}
%ay''+by'+cy=0,\ a,b,c\in\R,\ a\neq 0.
%\end{equation} 
% }
%\bigskip
%
%\uncover<2->{Para resolver uma equação do tipo (\ref{coef_const}) vamos nos inspirar no caso de 1\fm ordem. Uma EDO linear homogêna de 1\fm com coeficientes constantes é da forma
%$$ay'+by=0,\ a,b\in\R,\ a\neq 0.$$
%Sabemos que as soluções para esta equação são $\dps y(t)=ce^{bt/a}$. Neste caso é natural supor que uma solução da EDO (\ref{coef_const}) seja da forma $y(t)=e^{\lambda t}$ para alguma constante $\lambda$. Daí, substituindo em (\ref{coef_const}) temos que
%$$a\lambda^2e^{\lambda t}+b\lambda e^{\lambda t}+ce^{\lambda t}=0\Leftrightarrow a\lambda^2+b\lambda+c=0.$$
%A última equação é dita \dt{equação característica.}}
%
%\uncover<3->{\begin{exe} Determinar as soluções da equação:
%\begin{enumerate}[a)]
%\item $y''+y'-2y=0$
%\item $y''+2y'+y=0$
%\item $y''-2y+2=0$
%\end{enumerate}
%\end{exe}}

%
%\end{scriptsize}
%\end{frame}

\subsection*{EDO's Não-Homogêneas}



\begin{frame}
\frametitle{Equações não-homogêneas }
%\begin{scriptsize}

É fácil ver que se {\color{red}$y_p(t)$} é uma solução de uma EDO não-homogênea
\[y''+{\color{blue}p(t)}y'+{\color{blue}q(t)}y={\color{blue}f(t)},\]
 {\color{cyan}$y_1$} e {\color{cyan}$y_2$} são soluções fundamentais da EDO homogênea correspondente, então a solução geral da equação não-homogênea é
\[y(t)={\color{red}y_p(t)}+c_1{\color{cyan}y_1(t)}+c_2{\color{cyan}y_2(t)}.\]



\end{frame}

\begin{frame}
\begin{exe}
Se no modelo de vibrações mecânicas existir uma força externa $F=F(t)$, então teremos um problema de \dt{oscilação forçada} que é modelado pela  seguinte equação não-homogênea:
\[{\color{blue}m}y''+{\color{orange}\gamma} y'+{\color{red}k}y =F(t).\]

Neste caso, como determinar a solução do  seguinte problema  cuja força externa é periódica?
\[y''+2y'+y=2\cos(t).\]
\end{exe}
\end{frame}

\begin{frame}
\begin{center}
\includegraphics[scale=0.5]{osc-for-ex1.png}
\end{center}
\end{frame}




\begin{frame}
\frametitle{Método dos Coeficientes a Determinar }

\uncover<1->{Este método funciona para qualquer EDO não-homogênea com coeficientes constantes
\[ay''+by'+cy=F_1(t)+F_2(t)+\cdots+F_k(t),\]
onde  
\[F_i(t)=e^{{\color{blue}\alpha }t}[(a_0+\ldots+a_nt^{{\color{teal}n}}) \cos({\color{red}\beta}t)+(b_0+\ldots+b_mt^{{\color{teal}m}})\sen({\color{red}\beta} t)].\]

Neste caso, deve-se procurar, para cada $F_i$, uma solução particular da forma
\[y_p(t)=t^{{\color{cyan}s}}e^{{\color{blue}\alpha } t}[(A_0+\ldots+A_qt^{\color{teal}q})\cos({\color{red}\beta} t)+(B_0+\ldots+B_qt^{{\color{teal}q}})\sen({\color{red}\beta} t)],\]
em que ${\color{teal}q}=\max\{{\color{teal}n},{\color{teal}m}\}$,  ${\color{cyan}s}$ é o menor inteiro não-negativo que garante que  nenhuma parcela de $y_p$ seja solução da equação homogênea correspondente e $A_0,\ldots,A_q,B_0,\ldots,B_q$ são coeficientes a serem determinados. 



}

\end{frame}

\begin{frame}
\begin{exe} Encontre a solução geral das seguintes equações:
\begin{enumerate}[a]
\item $y''+y'=2+t^2$.
\item $y''-2y'+y=e^{t}+t$
\item $y''+4y=e^t\cos t$

\end{enumerate}
\end{exe}
\end{frame}


\begin{frame}{Vibrações Mecânicas Forçadas}
conteúdo...
\end{frame}



\begin{frame}
\frametitle{ Método da Variação dos Parâmetros}
%\begin{scriptsize}

Este método funciona para qualquer EDO linear de 2\fm ordem 
\[y''+{\color{blue}p(t)}y'+{\color{blue}q(t)}y={\color{blue}f(t)},\]
para o qual se conheça duas soluções fundamentais da equação homogênea correspondente em um intervalo $I$ onde o wronskiano é não nulo.
\bigskip

Sabemos que a solução geral da equação homogênea correspondente é
$$y(t)={\color{red}c_1}y_1(t)+{\color{red}c_2}y_2(t).$$
O método da variação dos parâmetros consiste em procurar uma solução particular da EDO não homogênea que tenha  a forma da solução geral da homogênea, mas substituindo os parâmetros $c_1$ e $c_2$ por funções a determinar $u_1(t)$ e $u_2(t)$, respectivamente, ou seja, da forma
\[y(t)={\color{red}u_1(t)}y_1(t)+{\color{red}u_2(t)}y_2(t),\]
com a condição de que
\[{\color{red}u'_1(t)}y_1(t)+{\color{red}u'_2(t)}y_2(t)=0\]

%\end{scriptsize}
\end{frame}

\begin{frame}{Fórmula Geral}
Resolvendo-se o sistema anterior, chega-se a:
\[{\color{red}u_1(t)=-\int \frac{y_2(t)f(t)}{W[y_1,y_2](t)}\,dt,\ \text{ e } u_2(t)=\int \frac{y_1(t)f(t)}{W[y_1,y_2](t)}\,dt}\]
\begin{exe}
 Encontre a solução do PVI
 $$\left\{
 \begin{array}{l}
 y''+y=\sec t\\
 \\
 y(0)=1,\ y'(0)=-2.
 \end{array}\right.$$
 \end{exe} 
\end{frame}


\begin{frame}
\frametitle{ }
\begin{exe} 

\begin{enumerate}
\item Determine a solução geral das EDO 
\[y''+2y'+2y=\frac{e^{-x}}{\cos^3(x)}.\]
\item Verifique que $y_1(t)=t^2$ e $y_2(t)=t^{-1}$ são soluções fundamentais da EDO no intervalo $(0,\infty)$
\[t^2y''-2y=3t^2-1\]
e determine a soução geral.
\end{enumerate}
\end{exe}
%\end{scriptsize}
\end{frame}



%
%
%\begin{frame}
%\frametitle{ }
%
%
%\center{\Huge{FIM !!!}
%\bigskip
%
%\includegraphics[scale=0.4]{ming.eps}}
%
%\end{frame}







\end{document}



